\let\negmedspace\undefined
\let\negthickspace\undefined
\documentclass[journal]{IEEEtran}
\usepackage[a5paper, margin=10mm, onecolumn]{geometry}
%\usepackage{lmodern} % Ensure lmodern is loaded for pdflatex
\usepackage{tfrupee} % Include tfrupee package

\setlength{\headheight}{1cm} % Set the height of the header box
\setlength{\headsep}{0mm}     % Set the distance between the header box and the top of the text

\usepackage{gvv-book}
\usepackage{gvv}
\usepackage{cite}
\usepackage{amsmath,amssymb,amsfonts,amsthm}
\usepackage{algorithmic}
\usepackage{graphicx}
\usepackage{textcomp}
\usepackage{xcolor}
\usepackage{txfonts}
\usepackage{listings}
\usepackage{enumitem}
\usepackage{mathtools}
\usepackage{gensymb}
\usepackage{comment}
\usepackage[breaklinks=true]{hyperref}
\usepackage{tkz-euclide} 
\usepackage{listings}
\def\inputGnumericTable{}                                 
\usepackage[latin1]{inputenc}                                
\usepackage{color}                                            
\usepackage{array}                                            
\usepackage{longtable}                                       
\usepackage{calc}                                             
\usepackage{multirow}                                         
\usepackage{hhline}                                           
\usepackage{ifthen}                                           
\usepackage{lscape}
\begin{document}

\bibliographystyle{IEEEtran}


\title{10.7.21}
\author{AI25BTECH11012 - GARIGE UNNATHI}
% \maketitle
% \newpage
% \bigskip
{\let\newpage\relax\maketitle}


\renewcommand{\thefigure}{\theenumi}
\renewcommand{\thetable}{\theenumi}
\setlength{\intextsep}{10pt} % Space between text and floats


\numberwithin{equation}{enumi}
\numberwithin{figure}{enumi}

\vspace{-1cm}

\textbf{Question}:\\
Find the points on the curve $x^2$ + $y^2$ - 2x - 3 = 0 at which the tangents are parallel
to the x-axis.


\textbf{Solution: }\\

For the given curve 
\begin{align}
    \vec{u} = \myvec{-1\\0}\\
    f = -3
\end{align}

Hence the centre and radius are given as :

\begin{align}
 \vec{c} = -\vec{u} = \myvec{1\\0}\\
 r = \sqrt{\lVert\vec{u}\rVert^2 - f} = 2
\end{align}

For a circle the points of contact are given by the formula :
\begin{align}
\vec{q}_{ij} = \left (\pm r \frac{\vec{n}_j}{{\lVert\vec{n}_j\rVert}} - \vec{u}\right)
\end{align}
Where i,j = 1,2

Since tangents are parallel to x - axis ,the normal is given as 
\begin{align}
  \vec{n} = \myvec{0\\1}
\end{align}

Solving 0.5 we get :

\begin{align}
 \vec{q} = (\pm 2\myvec{0\\1} - \myvec{-1\\0})\\
= \myvec{1\\-2} , \myvec{1\\2}
\end{align}


Hence the two points of contact are 
\begin{align}
    \myvec{1\\2} \quad and \quad \myvec{1\\-2}
\end{align}



\begin{figure}[h!]
   \centering
   \includegraphics[width=0.7\linewidth]{/Users/unnathi/Documents/ee1030-2025/ai25btech11012/matgeo/10.7.21/figs/fig.png}
   \caption{}
   \label{stemplot}
\end{figure}



\end{document}


